\section*{\texorpdfstring{\S\CurrentExerciseGroup}{第{\CurrentExerciseGroup}节}}
\phantomsection
\addcontentsline{toc}{section}{第{\CurrentExerciseGroup}节习题}

\begin{enumerate}
\def\labelenumi{\arabic{enumi})}
\item
  设 X 为局部紧空间 T 的局部紧子空间。证明：对 T 上的每个测度 \(\lambda\)，测度 \(\lambda_{X}\) 的支集是测度 \(\varphi_{X} \cdot \lambda\) 的支集在 X 上的迹。
\item
  设 \(\mathrm{T}\) 为局部紧空间，\(\mathrm{X}\) 为 \(\mathrm{T}\) 的局部紧子空间，\(\mu\) 为 \(\mathrm{T}\) 上的正测度；令 \(\nu = \varphi_{\mathrm{X}} \cdot \mu = j(\mu_{\mathrm{X}})\)（\(j\) 为 \(\mathrm{X}\) 到 \(\mathrm{T}\) 中的典则单射）。
\end{enumerate}

\begin{enumerate}
\def\labelenumi{\alph{enumi})}
\item
  证明：每个 \(\mu_{\mathrm{X}}\) -可忽略的 \(\mathbf{X}\) 的子集都是 \(\nu\) -可忽略的。
\item
  对每个数值函数 \(g \geqslant 0\)（定义在 \(X\) 上），证明 \(\int^{*} g \, d\mu_{X} = \int_{X}^{*} g \, d\nu\)（利用 \(a\) ) 与 §3 的公式 (7))。
\item
  设 \(\mathbf{g}\) 为定义在 \(X\) 上的函数，取值于 \(\overline{\mathbf{R}}\) 或一个 Banach 空间，\(\mathbf{g}'\) 为 \(\mathbf{g}\) 到 \(T\) 上的延拓，它在 \(T - X\) 上等于 0。为使 \(\mathbf{g}\) 是 \(\mu_{X}\) -可积的，必须且只需 \(\mathbf{g}'\) 是 \(\nu\) -可积的，此时 \(\int \mathbf{g} d\mu_X = \int_X \mathbf{g} d\nu\)（利用 \(b\))）。
\end{enumerate}

\begin{enumerate}
\def\labelenumi{\arabic{enumi})}
\setcounter{enumi}{2}
\tightlist
\item
  记号与习题 2 中的相同，假设 \(\mathbf{X}\) 是 \(\mathbf{T}\) 的开子集。
\end{enumerate}

\begin{enumerate}
\def\labelenumi{\alph{enumi})}
\item
  对每个数值函数 \(g \geqslant 0\)（定义在 \(X\) 上），证明 \(\int^{*} g d\mu_{X} = \int_{X}^{*} g d\mu\)（利用习题 2、§5 的命题 3 与 §1 的命题 4）。
\item
  设 g 为定义在 X 上的函数，取值于 \(\overline{R}\) 或一个 Banach 空间，\(g'\) 为 g 到 T 上的延拓，它在 T−X 上等于 0。为使 g 是 \(\mu_{X}\) -可积的，必须且只需 \(g'\) 是 \(\mu\) -可积的，此时 \(\int g d\mu_{X} = \int_{X} g d\mu\)（利用 a)）。
\end{enumerate}

\begin{enumerate}
\def\labelenumi{\arabic{enumi})}
\setcounter{enumi}{3}
\item
  记号同习题 2，证明：若 X 在 T 中闭，则对每个定义在 T 上的数值函数 \(f \geqslant 0\)，\(\int^{*}(f \circ j) d\mu_{X} = \int^{*} f d\nu\)（注意映射 \(j\) 是逆紧的，并利用 §4 的命题 2）。在同样条件下，若 \(\mathbf{f}\) 是 T 到 \(\overline{\mathbf{R}}\) 或到一个 Banach 空间中的映射，则为使 \(\mathbf{f} \circ j\) 是 \(\mu_{X}\) -可积的，必须且只需 \(\mathbf{f}\) 是 \(\nu\) -可积的，此时 \(\int(\mathbf{f} \circ j)\mu_{X} = \int \mathbf{f} d\nu\)（§4, 定理 2）。
\item
  设 \(\mathbf{T}\) 与 \(\mu\) 为 Ch. IV, §1 的习题 5 中定义的局部紧空间与测度，并设 \(\mathbf{D}\) 为 \(\mathbf{T}\) 中形如 \((0,y)\) 的点所成的（闭）集。
\end{enumerate}

\begin{enumerate}
\def\labelenumi{\alph{enumi})}
\item
  证明：测度 \(\mu_{\mathrm{D}}\)（由 \(\mu\) 诱导到 D 上的）为零；然而 \(\int^{*}\varphi_{\mathrm{D}}d\mu = +\infty\)（参见习题 3 a)）。
\item
  设 \(\mathrm{X} = \mathrm{T} - \mathrm{D}\)，并设 \(j\) 为 \(\mathrm{X}\) 到 \(\mathrm{T}\) 中的典则单射；则 \(j(\mu_{\mathrm{X}}) = \mu\)，但 \(\int^{*} (\varphi_{\mathrm{D}} \circ j) d\mu_{\mathrm{X}} = 0\)，\(\int^{*} \varphi_{\mathrm{D}} d\mu = +\infty\)（参见习题 4）。
\end{enumerate}

\begin{enumerate}
\def\labelenumi{\arabic{enumi})}
\setcounter{enumi}{5}
\item
  设 T 为局部紧空间，X 为 T 的局部紧子空间，j 为 X 到 T 中的典则单射。设 \(\lambda\) 为 X 上的正测度，使对每个紧集 \(K \subset T\)，\(K \cap X\) 是本质 \(\lambda\) -可积的。证明：测度 \(\nu = j(\lambda)\) 是 T 上最小的正测度 \(\rho\)（使 \(\rho_{X} = \lambda\)），且其支集是 \(\lambda\) 的支集在 T 中的闭包。
\item
  设 X、Y 为两个局部紧空间，\(\pi: X \to Y\) 为连续映射，T 为 Y 的局部紧子空间，\(S = \overline{\pi}^{-1}(T)\)，\(\pi_{T}: S \to T\) 为在 S 上与 \(\pi\) 重合的映射，\(\nu\) 为 Y 上的正测度。证明下列三个条件等价：
\end{enumerate}

\begin{enumerate}
\def\labelenumi{\alph{enumi})}
\item
  若 \(\nu_{\mathrm{T}}\) 是 \(\nu\) 在 \(\mathbf{T}\) 上诱导的测度，则存在一个正测度 \(\mu_{\mathrm{T}}\)（在 \(\mathbf{S}\) 上），使 \(\pi_{\mathrm{T}}(\mu_{\mathrm{T}}) = \nu_{\mathrm{T}}\)。
\item
  存在一个正测度 \(\lambda\)（在 \(X\) 上）使 \(\pi (\lambda) = \nu_{\mathrm{T}}\)。
\item
  存在 S 上的正测度 \(\lambda_{\mathrm{S}}\) 使 \((\pi|\mathrm{S})(\lambda_{\mathrm{S}})=\nu_{\mathrm{T}}\)。
\end{enumerate}

\begin{enumerate}
\def\labelenumi{\arabic{enumi})}
\setcounter{enumi}{7}
\item
  设 X、Y 为两个局部紧空间，\(\pi: X \to Y\) 为连续映射，\(\nu\) 为 Y 上集中在 \(\pi(X)\) 上的正测度。假设对每个 \(y \in \pi(X)\)，存在 y 在 Y 中的一个邻域 V 及一个正测度 \(\lambda_{\mathrm{V}}\)（在 \(\overline{\pi}^{-1}(V)\) 上），使 \(\lambda_{V}\) 在映射 \(\pi_{\mathrm{V}}: \overline{\pi}^{-1}(V) \to V\)（它与 \(\pi\) 重合）下的像等于诱导测度 \(\nu_{V}\)。证明：在这些条件下，存在 X 上的一个正测度 \(\mu\) 使 \(\pi(\mu) = \nu\)。（考察偶 \((G, \lambda)\) 所成之集 G，其中 G 在 Y 中开，而 \(\lambda\) 是 X 上使 \(\pi(\lambda) = \nu_{\mathrm{G}}\) 的正测度。用关系 \(\langle G_{1} \subset G_{2} \text{ 且 } \lambda_{1} \leqslant \lambda_{2} \rangle\)（定义在 \((G_{1}, \lambda_{1})\) 与 \((G_{2}, \lambda_{2})\) 之间）给 G 排序。先证明有序集 G 是归纳的，然后借助 Zorn 引理与习题 7 完成论证。）
\item
  设 X、Y 为两个局部紧空间，\(\pi: X \to Y\) 为连续映射。称映射 \(\pi\) 为保守映射，如果对每个正测度 \(\nu\)（Y 上的、集中在 \(\pi(X)\) 上的），存在 X 上的一个正测度 \(\mu\) 使 \(\nu = \pi(\mu)\)。
\end{enumerate}

\begin{enumerate}
\def\labelenumi{\alph{enumi})}
\item
  设 X、Y、Z 为三个局部紧空间，\(\pi:X\to Y\) 与 \(\pi^{\prime}:Y\to Z\) 为两个连续映射；令 \(\pi^{\prime\prime}=\pi^{\prime}\circ\pi\)，并假设映射 \(\pi\) 是满射。证明：若 \(\pi\) 与 \(\pi^{\prime}\) 是保守的，则 \(\pi^{\prime\prime}\) 也是；反之，若 \(\pi^{\prime\prime}\) 是保守的，则 \(\pi^{\prime}\) 也是。
\item
  假设对每个 \(y \in \pi(X)\)，存在 y 在 Y 中的一个开邻域 V，使得若 \(\pi_{\mathrm{V}} : \overline{\pi}^{-1}(\mathrm{V}) \to \mathrm{V}\) 是与 \(\pi\) 重合的映射，则存在一个与 \(\pi_{V}\) 相联的连续截口（GT, I, §3, No.~5）。证明 \(\pi\) 是保守的（利用习题 8）。
\item
  设 Y 为 R 的区间 \([0,1]\)，X 为赋有离散拓扑的集合 Y，\(\pi:X\to Y\) 为恒等映射。证明映射 \(\pi\) 不是保守的。
\end{enumerate}

\begin{enumerate}
\def\labelenumi{\arabic{enumi})}
\setcounter{enumi}{9}
\item
  设 \(X, Y\) 为两个局部紧空间。证明：每个逆紧连续映射 \(\pi\)（把 \(X\) 映入 \(Y\) 中）都是保守的。（该假设蕴含：映射 \(f \mapsto f \circ \pi\)（从 \(\mathcal{K}(Y)\) 到 \(\mathcal{C}(X)\) 中）的像 \(E\) 含于 \(\mathcal{K}(X)\) 中；然后，对每个正测度 \(\nu\)（在 \(Y\) 上），把 TVS, II, §3, No.~1 的命题 1 的推论应用于子空间 \(E\)（\(\mathcal{K}(X)\) 的）以及正线性形式 \(f \circ \pi \mapsto \nu(f)\)（\(E\) 上的）。）
\item
  设 X、Y 为两个局部紧空间，\(\pi: X \to Y\) 为连续映射，M 为 Y 的子集，使 \(\overline{\pi}^{-1}(M)\) 含于紧集的可数并中。则对每个集中在 M 上的 Y 上正测度 \(\nu\)，存在 X 上的正测度 \(\mu\) 使 \(\pi(\mu) = \nu\)（利用习题 10）。
\end{enumerate}
% semantic-fragments: legacy-input-seam
\relax{}
\setcounter{exercisegroup}{8}
\edef\CurrentExerciseGroup{\arabic{exercisegroup}}
